% ============================================================
%  Multimodal Volatility Forecasting — SSRN Preprint Draft
%  Author: Rishi Praval
%  Status: COMPLETE
% ============================================================

\documentclass[12pt, a4paper]{article}

% ---------- Packages ----------
\usepackage[margin=1in]{geometry}
\usepackage{amsmath, amssymb, amsthm}
\usepackage{graphicx}
\usepackage{booktabs}
\usepackage{hyperref}
\usepackage{xcolor}
\usepackage{setspace}
\usepackage{natbib}
\usepackage{caption}
\usepackage{subcaption}
\usepackage{microtype}
\usepackage{enumitem}
\usepackage{float}
\usepackage{tikz}
\usepackage{algorithm}
\usepackage{algpseudocode}

\onehalfspacing

\hypersetup{
    colorlinks=true,
    linkcolor=blue!70!black,
    citecolor=blue!70!black,
    urlcolor=blue!70!black
}

% ---------- Custom commands ----------
\newcommand{\hatσ}{\hat{\sigma}}

% ============================================================
\begin{document}

% ---------- Title ----------
\title{\textbf{Multimodal Volatility Forecasting via Gated Multi-Task Fusion\\
of Price, Text, Graph, and Macroeconomic Signals}}

\author{
    Rishi Praval \\
    \small Department of Computer Science (AI-ML), Adani University \\
    \small Ahmedabad, Gujarat, India \\
    \small \href{mailto:rishipraval007@gmail.com}{rishipraval007@gmail.com} \quad
    \href{https://github.com/imrishi007}{github.com/imrishi007}
}

\date{\today \\ \small \textit{Preprint — submitted to SSRN}}

\maketitle
\thispagestyle{empty}

% ---------- Abstract ----------
\begin{abstract}
Forecasting realized volatility in equity markets is a fundamental challenge in quantitative finance, with direct implications for options pricing, portfolio risk management, and dynamic hedging. Classical approaches such as HAR-RV exploit the well-documented long-memory structure of volatility but rely exclusively on past price history, ignoring the substantial information embedded in corporate disclosures, inter-firm dependencies, and macroeconomic regimes. We propose a gated multi-task fusion architecture that jointly encodes four heterogeneous signal modalities: price dynamics via a CNN-BiLSTM, fundamental text signals via FinBERT applied to SEC 10-K filings, inter-stock relational structure via a two-layer Graph Attention Network operating on supply-chain and competitive linkages, and macroeconomic context via a feed-forward MLP. A HAR-RV skip connection encodes a strong autoregressive prior, allowing the neural component to learn residual corrections over the classical benchmark rather than fitting volatility from scratch. Evaluated on 30 NASDAQ technology equities over a held-out 2024--2025 test period, our model achieves a volatility $R^2$ of 0.921 and a directional AUC of 0.591, closing 95.7\% of the benchmark gap relative to HAR-RV. These results suggest that multimodal fusion with appropriate inductive biases substantially advances the state of equity volatility forecasting while confirming the well-known asymmetry between the forecastability of second-moment and first-moment quantities in efficient markets.
\end{abstract}

\textbf{Keywords:} volatility forecasting, multimodal learning, graph attention networks, FinBERT, HAR-RV, multi-task learning, QLIKE loss

\newpage
\tableofcontents
\newpage

% ============================================================
\section{Introduction}
\label{sec:intro}
% ============================================================

Equity volatility forecasting occupies a central position in quantitative finance. Accurate volatility estimates are essential inputs to options pricing \citep{black1973pricing}, portfolio risk management \citep{markowitz1952portfolio}, and dynamic hedging strategies. Despite decades of research, producing reliable out-of-sample volatility forecasts remains an open challenge, particularly over multi-week horizons where the signal-to-noise ratio deteriorates sharply.

The dominant classical approach, the Heterogeneous Autoregressive model for Realized Volatility (HAR-RV), introduced by \citet{corsi2009simple}, decomposes realized volatility into daily, weekly, and monthly autoregressive components. HAR-RV is remarkably parsimonious and competitive out-of-sample, often outperforming far more complex models. However, it operates exclusively on past price history, ignoring the substantial information embedded in corporate filings, inter-firm economic relationships, and macroeconomic regime indicators.

Recent advances in deep learning and natural language processing have opened new avenues for financial forecasting. \citet{gu2020empirical} demonstrated that machine learning methods can systematically exploit nonlinearities in the cross-section of expected returns, while \citet{fischer2018deep} showed that LSTM networks capture temporal dependencies in asset returns that elude simpler models. These results motivate the hypothesis that deep learning architectures, when properly constrained by domain-appropriate inductive biases, can augment classical volatility models with information from non-traditional data sources.

In this work, we propose a gated multi-task fusion architecture that integrates four heterogeneous signal modalities: (1) historical price dynamics encoded by a CNN-BiLSTM, (2) fundamental textual signals from annual SEC 10-K filings processed by FinBERT \citep{yang2020finbert}, (3) inter-stock relational structure modeled by a Graph Attention Network (GAT) \citep{velivckovic2018graph}, and (4) macroeconomic context injected through a feed-forward MLP. Crucially, we introduce a HAR-RV skip connection that encodes the strong autoregressive prior directly into the architecture, allowing the model to learn residual corrections over the classical benchmark rather than fitting volatility from scratch.

Evaluated on 30 NASDAQ technology equities over a strictly held-out 2024--2025 test period, our model achieves a volatility $R^2$ of \textbf{0.921} and a directional AUC of \textbf{0.591}, closing \textbf{95.7\%} of the benchmark gap over HAR-RV. These results suggest that multimodal fusion, when combined with appropriate inductive biases, substantially advances the state of equity volatility forecasting.

The remainder of this paper is organized as follows. Section~\ref{sec:related} reviews related work. Section~\ref{sec:data} describes the dataset. Section~\ref{sec:method} presents the architecture. Section~\ref{sec:experiments} reports experimental results. Section~\ref{sec:discussion} discusses implications and limitations. Section~\ref{sec:conclusion} concludes.


% ============================================================
\section{Related Work}
\label{sec:related}
% ============================================================

\subsection{Classical Volatility Models}
The ARCH family of models, introduced by \citet{engle1982autoregressive}, established the foundational framework for modeling time-varying variance in financial time series. \citet{bollerslev1986generalized} generalized ARCH to the GARCH specification, which remains widely used in practice for option-implied and realized volatility estimation. \citet{corsi2009simple} proposed HAR-RV, which approximates long-memory behavior in realized volatility through a cascading linear structure of daily, weekly, and monthly components. Despite its simplicity, HAR-RV has proven remarkably competitive in out-of-sample forecasting benchmarks, routinely outperforming GARCH and stochastic volatility alternatives \citep{patton2011volatility}.

\subsection{Deep Learning for Volatility and Return Forecasting}
\citet{fischer2018deep} applied LSTM networks to stock return prediction across a broad cross-section of US equities, finding statistically significant improvements over random forests and logistic regression. \citet{gu2020empirical} conducted a comprehensive comparison of machine learning methods for asset pricing, concluding that neural networks with carefully chosen architectures and regularization achieve the best out-of-sample $R^2$ among 30,000+ individual stocks. Our work differs from these studies in two respects: first, we target realized volatility rather than expected returns; second, we explicitly integrate non-price information through a multimodal architecture rather than relying solely on price-derived features.

\subsection{NLP and Text in Finance}
The application of transformer-based language models to financial text has accelerated since the release of BERT. \citet{yang2020finbert} introduced FinBERT, a BERT model pre-trained on financial communications, demonstrating strong performance on sentiment classification of earnings call transcripts and analyst reports. Subsequent work has applied FinBERT to SEC filings for return prediction and event detection. However, the use of annual 10-K filings for \textit{volatility} forecasting specifically remains relatively unexplored. Our approach processes multiple sections of the 10-K filing (Item~1: Business, Item~1A: Risk Factors, Item~7: Management's Discussion and Analysis, Item~7A: Quantitative and Qualitative Disclosures, and Item~8: Financial Statements) to capture both forward-looking risk disclosures and fundamental operating characteristics.

\subsection{Graph Neural Networks in Finance}
Graph neural networks have been applied to stock prediction by encoding inter-firm relationships as graph structures. Most existing work constructs graphs from rolling correlation matrices or sector membership, which conflates statistical co-movement with economic linkage. \citet{velivckovic2018graph} introduced the GAT mechanism, which learns attention weights over a node's neighborhood. In our work, the inter-stock graph is constructed from hand-curated supply-chain, competitive, and operational linkages across seven relation types (supplier, platform competitor, advertising competitor, cloud competitor, consumer attention overlap, enterprise competitor, and index correlation), ensuring that graph edges carry semantic meaning beyond statistical artifact.

\subsection{Multi-Task Learning in Finance}
Joint prediction of volatility and return direction is relatively unexplored in the finance literature. Multi-task learning has been applied in NLP and computer vision to improve generalization by sharing representations across related objectives. In our setting, volatility and direction are economically related: periods of high volatility tend to coincide with directional regime changes. Our architecture exploits this by sharing a common trunk representation across volatility and direction heads, with a gating mechanism that learns to weight each modality differently for each task.


% ============================================================
\section{Data}
\label{sec:data}
% ============================================================

\subsection{Stock Universe}

We construct our study universe from 30 US technology-sector equities traded on NASDAQ. The initial ten stocks were selected from the largest-capitalization NASDAQ constituents, stratified across six technology sub-sectors (Table~\ref{tab:stocks}). The remaining twenty equities were added through a graph-motivated expansion procedure: since our architecture models inter-stock dependencies via a Graph Attention Network, we required each additional stock to have demonstrable operational or supply-chain linkages to at least one anchor stock. This ensures that graph edges carry semantic meaning beyond statistical return correlations.

\begin{table}[H]
\centering
\caption{Stock universe by sector}
\label{tab:stocks}
\begin{tabular}{ll}
\toprule
\textbf{Sector} & \textbf{Tickers} \\
\midrule
Core Technology (anchor) & AAPL, MSFT, GOOGL, AMZN, META, NVDA, AMD, INTC, TSLA, ORCL \\
Semiconductors           & QCOM, TXN, AVGO, MRVL, KLAC \\
Cloud / Software         & CRM, ADBE, NOW, SNOW, DDOG \\
Internet / Consumer      & NFLX, UBER, PYPL, SNAP \\
Hardware / Infrastructure& DELL, AMAT, LRCX \\
Diversified Technology   & IBM, CSCO, HPE \\
\bottomrule
\end{tabular}
\end{table}

Daily OHLCV price data and macroeconomic features were sourced from Yahoo Finance via the \texttt{yfinance} library. SEC 10-K filings were retrieved from the EDGAR database for each firm.

\subsection{Temporal Split}

Financial time series exhibit non-stationarity and regime dependence, making random cross-validation inappropriate. We adopt a strictly chronological train-validation-test split:

\begin{itemize}[noitemsep]
    \item \textbf{Training:} January 2016 -- December 2022 ($\approx$65\% of observations)
    \item \textbf{Validation:} January 2023 -- December 2023 ($\approx$10\%)
    \item \textbf{Test:} January 2024 -- December 2025 ($\approx$25\%)
\end{itemize}

No data from the validation or test periods is used during model training or hyperparameter selection, ensuring a fully realistic out-of-sample evaluation.

\subsection{Volatility Target}

The forecasting target is 60-day realized volatility, computed as the annualized rolling standard deviation of daily log-returns:

\begin{equation}
    \sigma^2_{i,t} = 252 \cdot \frac{1}{60} \sum_{k=1}^{60} \left( r_{i,t-k} - \bar{r}_{i,t} \right)^2
\end{equation}

where $r_{i,t} = \log(P_{i,t} / P_{i,t-1})$ is the daily log-return for stock $i$ at time $t$, and $\bar{r}_{i,t}$ is the mean log-return over the same 60-day window.

The 60-day horizon was selected after extensive ablation. Initial experiments at a 5-day horizon yielded near-random directional forecasts; the shift to 60 days represented the single largest performance improvement across all phases of development.

\subsection{Macroeconomic Features}

Twelve macroeconomic features are constructed to provide regime context (Table~\ref{tab:macro}). These capture market fear, yield curve structure, sector momentum, credit stress, and safe-haven demand --- collectively forming a compact representation of the macroeconomic environment at each forecast date.

\begin{table}[H]
\centering
\caption{Macroeconomic features}
\label{tab:macro}
\begin{tabular}{lll}
\toprule
\textbf{Feature} & \textbf{Source} & \textbf{Interpretation} \\
\midrule
VIX level              & \^{}VIX     & Market fear index \\
VIX daily change       & \^{}VIX     & Direction of fear \\
Long-term yields       & TLT         & 20-year Treasury ETF \\
Medium-term yields     & IEF         & 7--10 year Treasury ETF \\
Yield curve spread     & TLT $-$ IEF & Recession indicator \\
Broad market momentum  & SPY         & S\&P 500 direction \\
Tech sector momentum   & QQQ         & NASDAQ direction \\
Short-term rates       & SHY         & 1--3 year Treasury \\
Gold momentum          & GLD         & Safe-haven demand \\
High-yield credit      & HYG         & Credit stress \\
Investment-grade credit& LQD         & Credit quality \\
Market breadth         & Computed    & Cross-stock co-movement \\
\bottomrule
\end{tabular}
\end{table}


% ============================================================
\section{Methodology}
\label{sec:method}
% ============================================================

\subsection{Architecture Overview}

We propose a multimodal forecasting architecture that processes four classes of financial signals through dedicated encoders before integration via a learned gating mechanism. Figure~\ref{fig:arch} provides an overview.

\begin{figure}[H]
    \centering
    \begin{tikzpicture}[
        node distance=0.8cm and 1.2cm,
        block/.style={rectangle, draw, rounded corners, minimum width=2.0cm, minimum height=0.7cm, align=center, font=\small},
        encoder/.style={block, fill=blue!15},
        gate/.style={block, fill=orange!20},
        head/.style={block, fill=green!15},
        skip/.style={block, fill=red!10},
        arrow/.style={-stealth, thick},
    ]
        % Encoders
        \node[encoder] (price) {CNN-BiLSTM\\(256-d)};
        \node[encoder, right=of price] (text) {FinBERT\\(768-d)};
        \node[encoder, right=of text] (gat) {GAT\\(256-d)};
        \node[encoder, right=of gat] (macro) {Macro MLP\\(32-d)};

        % Projections
        \node[block, below=of price, fill=gray!10] (pp) {Proj 128-d};
        \node[block, below=of text, fill=gray!10] (tp) {Proj 128-d};
        \node[block, below=of gat, fill=gray!10] (gp) {Proj 128-d};
        \node[block, below=of macro, fill=gray!10] (mp) {Proj 128-d};

        % Gate
        \node[gate, below=2.0cm of tp, xshift=0.6cm] (gate) {Sigmoid Gate\\($\mathbf{g} \in [0,1]^4$)};

        % Trunk
        \node[block, below=of gate, fill=purple!10, minimum width=3cm] (trunk) {Shared Trunk\\(2-layer MLP, 256-d)};

        % Skip
        \node[skip, left=2.0cm of trunk] (skip) {HAR-RV\\Skip (32-d)};

        % Heads
        \node[head, below left=1.0cm and 0.3cm of trunk] (vol) {Vol Head\\(Softplus)};
        \node[head, below right=1.0cm and 0.3cm of trunk] (dir) {Dir Head\\(ListNet)};

        % Arrows
        \draw[arrow] (price) -- (pp);
        \draw[arrow] (text) -- (tp);
        \draw[arrow] (gat) -- (gp);
        \draw[arrow] (macro) -- (mp);

        \draw[arrow] (pp) -- (gate);
        \draw[arrow] (tp) -- (gate);
        \draw[arrow] (gp) -- (gate);
        \draw[arrow] (mp) -- (gate);

        \draw[arrow] (gate) -- (trunk);
        \draw[arrow] (skip) -- (trunk);
        \draw[arrow] (trunk) -- (vol);
        \draw[arrow] (trunk) -- (dir);
    \end{tikzpicture}
    \caption{Overview of the proposed multimodal gated fusion architecture. Four encoders process price, text, graph, and macro signals independently. Their 128-dimensional projections are weighted by a sigmoid gating network and fed through a shared 2-layer trunk. A HAR-RV skip connection injects lagged realized volatility features directly into the trunk, bypassing the gating mechanism. Two task-specific heads produce volatility predictions (Softplus activation) and directional rankings (ListNet loss).}
    \label{fig:arch}
\end{figure}

\subsection{Price Encoder: CNN-BiLSTM}

Daily OHLCV sequences of length 60 are processed by a hybrid CNN-BiLSTM encoder. Two one-dimensional convolutional layers (64 channels each, kernel size 3, stride 1, padding 1) with ReLU activations and batch normalization extract local price patterns and candlestick-level features across short windows. The resulting feature maps are passed to a 2-layer Bidirectional LSTM with hidden dimension 128 and dropout 0.2, which captures long-range temporal dependencies in both forward and backward directions. The final hidden state from the last time step is taken, yielding a fixed-dimensional price representation $\mathbf{z}_{\text{price}} \in \mathbb{R}^{256}$ (128 forward + 128 backward).

The input feature vector at each time step comprises 21 engineered features: raw OHLCV values, multi-horizon returns (1-, 5-, 10-day), rolling volatility (10-day), volume change, moving average ratios, Parkinson and Garman-Klass volatility estimators, 5-day realized volatility, volatility ratio, squared and absolute returns, a jump indicator, EWMA variance, and three HAR-RV autoregressive features (rv\_lag1d, rv\_lag5d, rv\_lag22d).

\subsection{Text Encoder: FinBERT on SEC 10-K Filings}

Fundamental information is extracted from annual SEC 10-K filings using FinBERT \citep{yang2020finbert}, specifically the \texttt{ProsusAI/finbert} pre-trained checkpoint. For each firm, the most recent 10-K filing available prior to each forecast date is retrieved to avoid lookahead bias. Five substantive sections of the filing are concatenated: Item~1 (Business), Item~1A (Risk Factors), Item~7 (Management's Discussion and Analysis of Financial Condition), Item~7A (Quantitative and Qualitative Disclosures About Market Risk), and Item~8 (Financial Statements and Supplementary Data). These sections collectively capture the firm's business model, forward-looking risk disclosures, management's narrative assessment of financial performance, and quantitative risk exposures.

The concatenated text is split into overlapping 512-token chunks (stride 128) and processed through FinBERT in FP16 precision for computational efficiency. The [CLS] token embedding from each chunk's final hidden layer is extracted, and the per-chunk embeddings are averaged to produce a single filing-level representation $\mathbf{z}_{\text{text}} \in \mathbb{R}^{768}$. An attention pooling layer is trained to weight chunk contributions by their informativeness.

\subsection{Relational Encoder: Graph Attention Network}

Inter-stock dependencies are modeled through a custom Graph Attention Network (GAT) implemented in pure PyTorch without external graph libraries. We construct a static graph $\mathcal{G} = (\mathcal{V}, \mathcal{E})$ where each node $v_i \in \mathcal{V}$ corresponds to a stock (30 nodes), and edges $\mathcal{E}$ are defined by seven semantically meaningful relation types:

\begin{enumerate}[noitemsep]
    \item \textbf{Supplier-to}: GPU/chip supply relationships (e.g., NVDA $\to$ TSLA, weight 0.95)
    \item \textbf{Platform competitor}: operating system and ecosystem competition (e.g., AAPL $\leftrightarrow$ GOOGL)
    \item \textbf{Advertising competitor}: digital advertising market competition (e.g., GOOGL $\leftrightarrow$ META)
    \item \textbf{Cloud competitor}: cloud infrastructure competition (e.g., AMZN $\leftrightarrow$ MSFT $\leftrightarrow$ GOOGL)
    \item \textbf{Consumer attention overlap}: competition for consumer mindshare (e.g., TSLA $\to$ AAPL)
    \item \textbf{Enterprise competitor}: enterprise software competition (e.g., ORCL $\to$ MSFT)
    \item \textbf{Index correlation}: co-movement within market indices
\end{enumerate}

Each edge carries a manually assigned prior weight reflecting the strength of the economic linkage. The graph is made bidirectional (excluding self-loops) and augmented with self-loops for all nodes.

Each node is initialized with the corresponding stock's CNN-BiLSTM representation, projected to the GAT input dimension of 256 (4 heads $\times$ 64 hidden units per head). The architecture comprises two GAT layers, each with 4 attention heads computing 64-dimensional representations, followed by ELU activation, residual connections, and LayerNorm. The attention mechanism follows \citet{velivckovic2018graph}:

\begin{equation}
    \mathbf{z}_{\text{graph},i} = \text{LayerNorm}\!\left(\mathbf{h}_i + \text{ELU}\!\left( \sum_{j \in \mathcal{N}(i)} \alpha_{ij} \mathbf{W} \mathbf{h}_j \right)\right)
\end{equation}

where $\alpha_{ij}$ are learned multi-head attention coefficients, and the residual connection ensures gradient flow through deep graph layers. The output dimensionality is $\mathbf{z}_{\text{graph}} \in \mathbb{R}^{256}$ (4 heads $\times$ 64).

\subsection{Macro Encoder: Feed-Forward MLP}

The twelve macroeconomic features (Table~\ref{tab:macro}) are concatenated and passed through a three-layer feed-forward MLP with architecture: Linear(12, 64) $\to$ LayerNorm $\to$ GELU $\to$ Dropout(0.2) $\to$ Linear(64, 64) $\to$ LayerNorm $\to$ GELU $\to$ Dropout(0.2) $\to$ Linear(64, 32). The macro encoder produces a context vector $\mathbf{z}_{\text{macro}} \in \mathbb{R}^{32}$. All macro features are normalized using statistics from the training period only, with lagged values (1-day lag) to prevent any form of lookahead bias.

\subsection{Gated Multi-Task Fusion}

The four encoder representations $\{\mathbf{z}_{\text{price}}, \mathbf{z}_{\text{text}}, \mathbf{z}_{\text{graph}}, \mathbf{z}_{\text{macro}}\}$ are first projected to a common 128-dimensional space via per-modality projectors (each consisting of a Linear layer, LayerNorm, GELU activation, and Dropout). The projected representations are integrated via a learned sigmoid gating mechanism. Gating weights $\mathbf{g} \in [0,1]^4$ are produced by a two-layer network that takes as input the concatenation of all projected modality representations and a 5-dimensional earnings surprise conditioning vector:

\begin{equation}
    \mathbf{g} = \text{Sigmoid}\!\left( \mathbf{W}_2 \, \text{GELU}\!\left(\mathbf{W}_1 \,[\mathbf{z}_{\text{price}}^{\text{proj}} \| \mathbf{z}_{\text{text}}^{\text{proj}} \| \mathbf{z}_{\text{graph}}^{\text{proj}} \| \mathbf{z}_{\text{macro}}^{\text{proj}} \| \mathbf{s}] + \mathbf{b}_1\right) + \mathbf{b}_2 \right)
\end{equation}

where $\mathbf{s} \in \mathbb{R}^5$ encodes earnings surprise features (surprise percentage, normalized surprise magnitude, recency decay, trailing 3-quarter beat rate, and EPS revision direction). A modality mask zeros out gates for unavailable modalities, and the remaining gates are renormalized to sum to one. This allows the model to adaptively suppress uninformative modalities on a per-sample basis:

\begin{equation}
    \mathbf{z}_{\text{fused}} = \sum_{k} g_k \cdot \mathbf{z}_k^{\text{proj}}
\end{equation}

The fused representation (512 dimensions) feeds into a shared 2-layer trunk with hidden dimension 256, LayerNorm, GELU activations, and dropout (0.3 for the first layer, 0.15 for the second), which then branches into two task-specific heads.

\subsection{HAR-RV Skip Connection}

A key design innovation is a direct skip connection that injects the three HAR-RV autoregressive features (rv\_lag1d, rv\_lag5d, rv\_lag22d) into the shared trunk, bypassing the gating mechanism entirely. The raw HAR-RV features are extracted from the last time step of each 60-day input window and projected through a small near-linear pathway:

\begin{equation}
    \mathbf{z}_{\text{HAR}} = \text{GELU}\!\left(\text{LayerNorm}\!\left(\mathbf{W}_{\text{HAR}} \cdot [\text{rv}_{t-1}, \bar{\text{rv}}_{t-5}, \bar{\text{rv}}_{t-22}]\right)\right) \in \mathbb{R}^{32}
\end{equation}

The trunk input is then the concatenation of the 512-dimensional gated fusion and the 32-dimensional HAR-RV skip, yielding 544 dimensions total. Crucially, the HAR-RV path is \textit{ungated}: it always contributes to the prediction, encoding the strong prior that lagged realized volatility is the single most informative predictor. The neural component thus learns to predict the \textit{residual} that the HAR-RV model cannot explain, rather than fitting volatility from scratch.

Empirically, this skip connection was responsible for the most significant single performance jump in our development cycle, improving volatility $R^2$ from 0.867 (without skip) to 0.921 (with skip).

\subsection{Loss Functions}

\paragraph{Volatility head --- QLIKE loss.}
Mean squared error treats over- and under-prediction symmetrically. In risk management contexts, however, underestimating volatility is substantially more costly than overestimating it --- an underpredicted VaR can result in capital insufficiency, while an overpredicted one is merely conservative. We therefore train the volatility head with the QLIKE loss \citep{patton2011volatility}, which asymmetrically penalizes underestimation:

\begin{equation}
    \mathcal{L}_{\text{QLIKE}} = \frac{1}{N} \sum_{i=1}^{N} \left( \frac{\sigma^2_{\text{true},i}}{\hat{\sigma}^2_i} + \log \hat{\sigma}^2_i \right)
\end{equation}

The volatility head uses a Softplus output activation to guarantee strictly positive predictions, with an $\epsilon=10^{-6}$ floor for numerical stability. QLIKE is also the standard evaluation metric in the volatility forecasting literature, making it appropriate both as a training objective and as a model selection criterion.

\paragraph{Direction head --- ListNet ranking loss.}
Rather than framing directional prediction as binary classification, we adopt a learning-to-rank formulation using ListNet \citep{cao2007learning}. At each time step, the model scores each stock and the ListNet loss trains it to reproduce the correct relative ranking of cross-sectional returns. The loss operates on the softmax distribution over predicted scores versus the softmax distribution over true returns:

\begin{equation}
    \mathcal{L}_{\text{ListNet}} = -\sum_{i} P_{\text{true}}(i) \log P_{\text{pred}}(i)
\end{equation}

where $P_{\text{true}}$ and $P_{\text{pred}}$ are softmax-normalized distributions over true and predicted scores respectively, with a temperature parameter $\tau = 5.0$ applied to the true labels to control distribution sharpness. This formulation is more appropriate than binary classification for cross-sectional equity selection, where relative performance matters more than absolute direction.

\paragraph{Combined loss.}
The total training objective is a weighted combination:

\begin{equation}
    \mathcal{L} = 0.85 \cdot \mathcal{L}_{\text{QLIKE}} + 0.15 \cdot \mathcal{L}_{\text{ListNet}}
\end{equation}

The 0.85/0.15 weighting reflects the volatility-primary design: the directional task serves as an auxiliary regularizer that encourages the shared trunk to encode directionally informative features, while the bulk of the optimization signal derives from the volatility objective.

\subsection{Training Details}

The model is trained for up to 60 epochs using AdamW \citep{loshchilov2019decoupled} with learning rate $5 \times 10^{-4}$, weight decay $10^{-3}$, and a cosine annealing schedule (minimum learning rate $10^{-6}$). Gradient norms are clipped at 1.0. Mixed-precision training (FP16) is used throughout via PyTorch's automatic mixed precision. Early stopping with patience 7 epochs is applied on the validation loss. Batch size is 4,096, which fits within 3.5 GB of GPU VRAM. A global random seed of 42 is used for reproducibility.

Monte Carlo (MC) Dropout is employed at inference time (30 forward passes) to estimate prediction uncertainty. The trunk dropout layer remains active during evaluation, with the mean prediction used for point estimates and the standard deviation used for uncertainty quantification and confidence-based filtering.


% ============================================================
\section{Experiments}
\label{sec:experiments}
% ============================================================

\subsection{Baselines}

We compare against the following baselines:
\begin{itemize}[noitemsep]
    \item \textbf{Historical Average (HA)}: Predicts the mean realized volatility from the training set. $R^2 = 0.348$.
    \item \textbf{HAR-RV} \citep{corsi2009simple}: The primary classical benchmark, using daily (1-day lagged), weekly (5-day average), and monthly (22-day average) realized volatility as predictors. $R^2 = 0.947$.
    \item \textbf{V2 Baseline}: Direction-primary multimodal fusion (Phase 11) using MSE loss for volatility. $R^2 = 0.335$.
    \item \textbf{Phase 12}: Volatility-primary multimodal fusion with QLIKE loss and MC Dropout, but without HAR-RV skip connection. $R^2 = 0.772$.
    \item \textbf{Phase 13}: Phase 12 architecture expanded to 30 stocks with HAR-RV features as CNN-BiLSTM inputs (not as skip connection). $R^2 = 0.867$.
\end{itemize}

\subsection{Development Trajectory}

Table~\ref{tab:phases} reports volatility $R^2$ and directional AUC across key development phases, illustrating the incremental contribution of each architectural innovation.

\begin{table}[H]
\centering
\caption{Performance across development phases on the held-out test set}
\label{tab:phases}
\begin{tabular}{llcc}
\toprule
\textbf{Phase} & \textbf{Description} & \textbf{Vol $R^2$} & \textbf{Dir AUC} \\
\midrule
Phase 6  & V1 Fusion (MSE, 10 stocks)         & 0.440 & 0.545 \\
Phase 11 & V2 Fusion + ListNet                & 0.335 & 0.565 \\
Phase 12 & Vol-Primary + QLIKE + MC Dropout   & 0.772 & 0.568 \\
Phase 13 & HAR-RV features + 30 stocks        & 0.867 & 0.585 \\
Phase 14 & HAR-RV Skip Connection (proposed)  & \textbf{0.921} & \textbf{0.591} \\
\midrule
HAR-RV   & Classical benchmark                & 0.947 & ---   \\
\bottomrule
\end{tabular}
\end{table}

Several observations are notable. First, switching from MSE to QLIKE loss (Phase 12) produced the largest single jump in volatility $R^2$ (+0.437), confirming that loss function design is a critical and underappreciated modeling choice in volatility forecasting. Second, expanding the universe from the initial 10 anchor stocks to 30 equities (Phase 13) improved both tasks, validating the graph-motivated stock selection procedure and demonstrating that the GAT benefits from a richer neighborhood structure. Third, the HAR-RV skip connection (Phase 14) contributed an additional 0.054 in $R^2$, closing 95.7\% of the remaining gap to the classical benchmark.

\subsection{Main Results}

Table~\ref{tab:main} presents the main results on the held-out test set with additive ablation, starting from the price-only encoder and progressively adding each modality.

\begin{table}[H]
\centering
\caption{Main results on held-out test set (2024--2025). Additive ablation from price-only to full model.}
\label{tab:main}
\begin{tabular}{lccc}
\toprule
\textbf{Model} & \textbf{Vol $R^2$} & \textbf{Dir AUC} & \textbf{Sharpe} \\
\midrule
HAR-RV (benchmark)            & 0.947 & ---   & ---  \\
Historical Average            & 0.348 & ---   & ---  \\
\midrule
V2 Baseline (MSE loss)        & 0.335 & 0.565 & ---  \\
Phase 12 (QLIKE, MC Dropout)  & 0.772 & 0.568 & ---  \\
Phase 13 (+ 30 stocks)        & 0.867 & 0.585 & ---  \\
\textbf{Phase 14 (+ HAR-RV skip)} & \textbf{0.921} & \textbf{0.591} & \textbf{0.62} \\
\bottomrule
\end{tabular}
\end{table}

The proposed model achieves a volatility $R^2$ of 0.921, recovering 95.7\% of the HAR-RV benchmark's explained variance while additionally providing directional forecasts (AUC 0.591) and uncertainty estimates via MC Dropout.

\subsection{Confidence Filtering}

We evaluate whether restricting predictions to high-confidence subsets improves directional accuracy. Confidence is defined as the maximum softmax probability across the two direction classes (UP/DOWN). At each coverage threshold, we retain only the most confident predictions and evaluate AUC and accuracy on the retained subset.

At approximately 30\% coverage (retaining only the top 30\% most confident predictions), the directional AUC improves to approximately 0.61, up from 0.591 at full coverage. This confirms that the model exhibits meaningful calibration: when it is confident in its directional forecast, it tends to be correct more often. The per-ticker analysis reveals heterogeneous confidence patterns, with liquid large-cap stocks (AAPL, MSFT, GOOGL) generally showing higher AUC at filtered thresholds than smaller or more volatile names.

The confidence filtering mechanism has direct practical implications: a trading desk could selectively act only on high-confidence signals, accepting lower trade frequency in exchange for higher per-trade accuracy.

\subsection{Backtest Results}

We evaluate the economic significance of our model's predictions through a simulated long-short trading strategy. The strategy is constructed as follows:

\begin{itemize}[noitemsep]
    \item \textbf{Signal}: Model-predicted probabilities rank stocks within each cross-section.
    \item \textbf{Position sizing}: Long the top 5 highest-conviction stocks, short the bottom 5, equal-weighted within each leg.
    \item \textbf{Rebalancing}: Every 5 trading days.
    \item \textbf{Holding period}: 20 trading days.
    \item \textbf{Transaction costs}: Tested at 0bp, 5bp, 10bp, and 20bp per trade.
\end{itemize}

At 5 basis points of transaction costs, the direction-driven long-short strategy produces a Sharpe ratio of 0.62 over the 2024--2025 test period. A VIX-adjusted volatility strategy, which uses the difference between model-predicted volatility and VIX-implied volatility ($\text{VIX} \times \beta_{\text{stock}}$) as a mispricing signal with uncertainty-weighted positions, provides an alternative economic application of the model's volatility forecasts.

These backtest results are indicative only. The simulation assumes no market impact, unlimited short availability, and instantaneous execution. Real-world implementation would face additional friction from bid-ask spreads, borrowing costs, and liquidity constraints, particularly for smaller names such as SNAP and DDOG.


% ============================================================
\section{Discussion}
\label{sec:discussion}
% ============================================================

\subsection{What the Volatility $R^2$ Means}

An $R^2$ of 0.921 against realized volatility is a strong result for a 60-day horizon, particularly given that the test period (2024--2025) is fully held out and unseen during training. For context, in the comprehensive volatility forecasting evaluation of \citet{patton2011volatility}, the best models typically achieve $R^2$ values in the 0.40--0.70 range at daily or weekly horizons; our higher value reflects both the smoother nature of 60-day realized volatility and the strong autoregressive prior encoded by the HAR-RV skip connection. Importantly, the HAR-RV benchmark achieves $R^2 = 0.947$, indicating that the autoregressive structure of realized volatility accounts for the majority of explainable variance. Our model recovers 95.7\% of this benchmark performance while additionally leveraging textual, relational, and macroeconomic signals.

\subsection{Why Direction is Harder}

A directional AUC of 0.591 is statistically significant but modest in absolute terms. This is consistent with the semi-strong form of market efficiency \citep{fama1970efficient}: if return direction were highly predictable from public information, the predictability would be rapidly arbitraged away. Volatility, by contrast, is a second-moment quantity that does not carry the same self-defeating property --- knowing future volatility does not in itself enable a risk-free arbitrage. This asymmetry between volatility and direction forecastability is a well-documented empirical regularity and provides theoretical grounding for our observed results.

\subsection{The Role of Loss Function Design}

The transition from MSE to QLIKE loss yielded the single largest improvement among all architectural changes (Phase 11 $R^2 = 0.335$ to Phase 12 $R^2 = 0.772$, a +0.437 improvement). This result underscores a point that is insufficiently appreciated in the applied deep learning literature: for domain-specific forecasting tasks, the choice of loss function can dominate the impact of architectural complexity. QLIKE's asymmetric penalty structure aligns with the economic asymmetry of volatility estimation errors, and its scale-sensitivity prevents the model from collapsing to mean predictions during high-volatility regimes.

\subsection{Limitations}

Several limitations should be noted. First, our universe is restricted to 30 large-cap US technology equities, limiting generalizability to other sectors, market capitalizations, or geographies. The graph structure and competitive linkages are specific to the technology sector and would require re-specification for other industries. Second, 10-K filings are annual documents; the text encoder cannot capture within-year information from quarterly earnings calls (10-Q), 8-K event filings, or analyst reports. Incorporating higher-frequency text sources is a natural extension. Third, the backtest assumes idealized execution and does not model market impact, which may be significant for less liquid names or larger portfolio sizes. Fourth, the 2024--2025 test period, while genuinely held out, represents a specific macroeconomic regime characterized by elevated interest rates and AI-driven sector rotation; performance may differ in other regimes. Finally, the graph structure is static and hand-curated. Dynamic graph construction from evolving supply-chain data or learned edge weights could improve the relational encoder's expressiveness.


% ============================================================
\section{Conclusion}
\label{sec:conclusion}
% ============================================================

We proposed a gated multi-task fusion architecture for joint volatility and directional forecasting of equities, integrating price, textual, relational, and macroeconomic signals through modality-specific encoders. The key architectural innovation --- a HAR-RV skip connection that grounds the model's prediction in the classical benchmark --- proved critical, contributing substantially to the final performance of $R^2 = 0.921$ on 60-day realized volatility.

More broadly, our results suggest that the dominant value of multimodal fusion in financial forecasting lies in volatility prediction rather than return direction --- a finding with direct implications for risk management and options market applications. We release our code and pre-trained model weights at \url{https://github.com/imrishi007/financial-document-analysis}.

Several avenues remain for future work. Expanding the stock universe beyond technology to other sectors and geographies would test the generalizability of the architecture. Incorporating higher-frequency text sources such as quarterly 10-Q filings, 8-K event filings, and earnings call transcripts could provide within-year textual signals that annual 10-K filings cannot capture. Intra-day price data and options-implied volatility surfaces offer additional signal sources for multi-horizon forecasting. Finally, replacing the static hand-curated graph with a dynamic, learned graph structure could enhance the relational encoder's adaptability to evolving inter-firm dependencies.


% ============================================================
% References
% ============================================================
\newpage
\bibliographystyle{plainnat}

\begin{thebibliography}{99}

\bibitem[Black and Scholes(1973)]{black1973pricing}
Black, F., \& Scholes, M. (1973).
The pricing of options and corporate liabilities.
\textit{Journal of Political Economy}, 81(3), 637--654.

\bibitem[Bollerslev(1986)]{bollerslev1986generalized}
Bollerslev, T. (1986).
Generalized autoregressive conditional heteroskedasticity.
\textit{Journal of Econometrics}, 31(3), 307--327.

\bibitem[Cao et al.(2007)]{cao2007learning}
Cao, Z., Qin, T., Liu, T.-Y., Tsai, M.-F., \& Li, H. (2007).
Learning to rank: From pairwise approach to listwise approach.
\textit{Proceedings of the 24th International Conference on Machine Learning (ICML)}, 129--136.

\bibitem[Corsi(2009)]{corsi2009simple}
Corsi, F. (2009).
A simple approximate long-memory model of realized volatility.
\textit{Journal of Financial Econometrics}, 7(2), 174--196.

\bibitem[Engle(1982)]{engle1982autoregressive}
Engle, R. F. (1982).
Autoregressive conditional heteroscedasticity with estimates of the variance of United Kingdom inflation.
\textit{Econometrica}, 50(4), 987--1007.

\bibitem[Fama(1970)]{fama1970efficient}
Fama, E. F. (1970).
Efficient capital markets: A review of theory and empirical work.
\textit{The Journal of Finance}, 25(2), 383--417.

\bibitem[Fischer and Krauss(2018)]{fischer2018deep}
Fischer, T., \& Krauss, C. (2018).
Deep learning with long short-term memory networks for financial market predictions.
\textit{European Journal of Operational Research}, 270(2), 654--669.

\bibitem[Gu, Kelly, and Xiu(2020)]{gu2020empirical}
Gu, S., Kelly, B., \& Xiu, D. (2020).
Empirical asset pricing via machine learning.
\textit{The Review of Financial Studies}, 33(5), 2223--2273.

\bibitem[Loshchilov and Hutter(2019)]{loshchilov2019decoupled}
Loshchilov, I., \& Hutter, F. (2019).
Decoupled weight decay regularization.
\textit{International Conference on Learning Representations (ICLR)}.

\bibitem[Markowitz(1952)]{markowitz1952portfolio}
Markowitz, H. (1952).
Portfolio selection.
\textit{The Journal of Finance}, 7(1), 77--91.

\bibitem[Patton(2011)]{patton2011volatility}
Patton, A. J. (2011).
Volatility forecast comparison using imperfect volatility proxies.
\textit{Journal of Econometrics}, 160(1), 246--256.

\bibitem[Veli\v{c}kovi\'{c} et al.(2018)]{velivckovic2018graph}
Veli\v{c}kovi\'{c}, P., Cucurull, G., Casanova, A., Romero, A., Li\`{o}, P., \& Bengio, Y. (2018).
Graph attention networks.
\textit{International Conference on Learning Representations (ICLR)}.

\bibitem[Yang et al.(2020)]{yang2020finbert}
Yang, Y., Uy, M. C. S., \& Huang, A. (2020).
FinBERT: A pretrained language model for financial communications.
\textit{arXiv preprint arXiv:2006.08097}.

\end{thebibliography}

\end{document}
